% GENERATED FILE. Edit canonical lessons, then run npm run generate:books.
% canonical-source-hash: fnv1a64:cb420e8d2679bf90
% canonical-lessons: LA-C129-ramus, LA-C129-truncus, LA-C129-semen, LA-C129-saxum, LA-C129-ripa

\chapter{Branch, Tree trunk, Seed, Rock, River bank}
\label{ch:la-branch-tree-trunk-seed-rock-river-bank}
\hlchaptermodality{129}

\begin{chapteropening}
\textbf{By the end of this chapter:} \emph{I can say a branch, a tree trunk, a seed, a rock and a river bank.}

Complete the last lesson of chapter 129: I can say a branch, a tree trunk, a seed, a rock and a river bank.
\end{chapteropening}

\section[rāmus, rāmī]{rāmus, rāmī — a branch}

\label{lesson:LA-C129-ramus}



\begin{quote}
Before the new one: say the Latin for a rose, then the Latin for a lily.
\end{quote}

\subsection*{You'll want to know: rāmus, rāmī}
\textbf{rāmus, rāmī} — \textquotedblleft{}a branch\textquotedblright{}.

\begin{cousinweb}
\textbf{rāmus} is probably related to \emph{rādīx}, \textquotedblleft{}root\textquotedblright{}, from an older \emph{*rādmus} (Lewis \& Short); both are traced to Proto-Indo-European \emph{*wreh2d-}, \textquotedblleft{}root, branch\textquotedblright{} (EtymDB 2.1).
\end{cousinweb}

\begin{grammarlens}[title={What you've built}]
One more word for this chapter.
\end{grammarlens}

\subsection*{Your turn}
\begin{itemize}
\raggedright
  \item \emph{Say it:} \emph{rāmus}
  \item \emph{Say it:} \emph{rāmus}, once more
  \item \emph{Say it:} say \emph{rāmus}
  \item \emph{Recall:} say the Latin for a rose, then the Latin for a lily, then say \emph{rāmus} again
\end{itemize}

\begin{tcolorbox}[breakable,colback=teal!4,colframe=teal!35!black,title={Before you move on}]
What does rāmus, rāmī mean? (\textquotedblleft{}A branch\textquotedblright{}.) Say it once more.
\end{tcolorbox}
\section[truncus, truncī]{truncus, truncī — a tree trunk}

\label{lesson:LA-C129-truncus}



\begin{quote}
Before the new one: say the Latin for a lily, then the Latin for a branch.
\end{quote}

\subsection*{You'll want to know: truncus, truncī}
\textbf{truncus, truncī} — \textquotedblleft{}a tree trunk\textquotedblright{}.

\begin{cousinweb}
\textbf{truncus} is the adjective \emph{truncus}, \textquotedblleft{}cut off, maimed\textquotedblright{} (Lewis \& Short), used as a noun: the stem of a tree, without its branches. English \emph{trunk} comes from it.
\end{cousinweb}

\begin{grammarlens}[title={What you've built}]
One more word for this chapter.
\end{grammarlens}

\subsection*{Your turn}
\begin{itemize}
\raggedright
  \item \emph{Say it:} \emph{truncus}
  \item \emph{Say it:} \emph{truncus}, once more
  \item \emph{Say it:} say \emph{truncus}
  \item \emph{Recall:} say the Latin for a lily, then the Latin for a branch, then say \emph{truncus} again
\end{itemize}

\begin{tcolorbox}[breakable,colback=teal!4,colframe=teal!35!black,title={Before you move on}]
What does truncus, truncī mean? (\textquotedblleft{}A tree trunk\textquotedblright{}.) Say it once more.
\end{tcolorbox}
\section[sēmen, sēminis]{sēmen, sēminis — a seed}

\label{lesson:LA-C129-semen}



\begin{quote}
Before the new one: say the Latin for a branch, then the Latin for a tree trunk.
\end{quote}

\subsection*{You'll want to know: sēmen, sēminis}
\textbf{sēmen, sēminis} — \textquotedblleft{}a seed\textquotedblright{}.

\begin{cousinweb}
\textbf{sēmen} is inherited, usually traced to Proto-Indo-European \emph{*seh1mn}, \textquotedblleft{}seed\textquotedblright{}, from \emph{*seh1-}, \textquotedblleft{}to sow\textquotedblright{} (EtymDB 2.1); \emph{serō}, \textquotedblleft{}to sow\textquotedblright{}, belongs to the same family (Lewis \& Short), and English \emph{seed} is a cousin.
\end{cousinweb}

\begin{grammarlens}[title={What you've built}]
One more word for this chapter.
\end{grammarlens}

\subsection*{Your turn}
\begin{itemize}
\raggedright
  \item \emph{Say it:} \emph{sēmen}
  \item \emph{Say it:} \emph{sēmen}, once more
  \item \emph{Say it:} say \emph{sēmen}
  \item \emph{Recall:} say the Latin for a branch, then the Latin for a tree trunk, then say \emph{sēmen} again
\end{itemize}

\begin{tcolorbox}[breakable,colback=teal!4,colframe=teal!35!black,title={Before you move on}]
What does sēmen, sēminis mean? (\textquotedblleft{}A seed\textquotedblright{}.) Say it once more.
\end{tcolorbox}
\section[saxum, saxī]{saxum, saxī — a rock, a boulder}

\label{lesson:LA-C129-saxum}



\begin{quote}
Before the new one: say the Latin for a tree trunk, then the Latin for a seed.
\end{quote}

\subsection*{You'll want to know: saxum, saxī}
\textbf{saxum, saxī} — \textquotedblleft{}a rock, a boulder\textquotedblright{}.

\begin{cousinweb}
\textbf{saxum} is often traced to Proto-Indo-European \emph{*sek-}, \textquotedblleft{}to cut\textquotedblright{} (EtymDB 2.1), the root of \emph{secāre}, \textquotedblleft{}to cut\textquotedblright{} (Lewis \& Short), but the link is not certain.
\end{cousinweb}

\begin{grammarlens}[title={What you've built}]
One more word for this chapter.
\end{grammarlens}

\subsection*{Your turn}
\begin{itemize}
\raggedright
  \item \emph{Say it:} \emph{saxum}
  \item \emph{Say it:} \emph{saxum}, once more
  \item \emph{Say it:} say \emph{saxum}
  \item \emph{Recall:} say the Latin for a tree trunk, then the Latin for a seed, then say \emph{saxum} again
\end{itemize}

\begin{tcolorbox}[breakable,colback=teal!4,colframe=teal!35!black,title={Before you move on}]
What does saxum, saxī mean? (\textquotedblleft{}A rock, a boulder\textquotedblright{}.) Say it once more.
\end{tcolorbox}
\section[rīpa, rīpae]{rīpa, rīpae — a river bank}

\label{lesson:LA-C129-ripa}



\begin{quote}
Before the new one: say the Latin for a seed, then the Latin for a rock.
\end{quote}

\subsection*{You'll want to know: rīpa, rīpae}
\textbf{rīpa, rīpae} — \textquotedblleft{}a river bank\textquotedblright{}.

\begin{cousinweb}
\textbf{rīpa} is probably inherited, from a Proto-Indo-European \emph{*h1reypeh2}, \textquotedblleft{}steep slope\textquotedblright{}, from a root for tearing (EtymDB 2.1); Greek \emph{ereipō}, \textquotedblleft{}to tear down\textquotedblright{}, is often compared. Lewis \& Short call it doubtful.
\end{cousinweb}

\begin{grammarlens}[title={What you've built}]
One more word for this chapter.
\end{grammarlens}

\subsection*{Your turn}
\begin{itemize}
\raggedright
  \item \emph{Say it:} \emph{rīpa}
  \item \emph{Say it:} \emph{rīpa}, once more
  \item \emph{Say it:} say \emph{rīpa}
  \item \emph{Recall:} say the Latin for a seed, then the Latin for a rock, then say \emph{rīpa} again
\end{itemize}

\begin{tcolorbox}[breakable,colback=teal!4,colframe=teal!35!black,title={Before you move on}]
What does rīpa, rīpae mean? (\textquotedblleft{}A river bank\textquotedblright{}.) Say it once more.
\end{tcolorbox}
